% s-311-p311-m002; original PDF 311, 312
\Aphorism{s-311-p311-m002}{365}
Sollte es nicht theoretische Einsichten geben, aus denen keine Regel für das practische Verhalten, keine normative Nutzanwendung herauszupressen wäre, bei denen die unanständige Beutegier des moralischen Menschen, der überall ein neues Du sollst! oder Du darfst nicht! auszuwittern sucht, einmal leer ausginge? Sollten nicht gerade die stärksten und allgemeinsten Erkenntnisse die ethisch unergiebigsten sein? — Bei allen bisherigen Versuchen, moralische Imperative an rein wissenschaftliche Indicative anzuschliessen, sprach immer wieder ein Irrthum das letzte Wort: man glaubte in aller Unschuld an die Allgemeingültigkeit von Normen und an die normative Anwendbarkeit allgemeingültiger Sätze, während es kaum einen schärferen Gegensatz und weiteren Abstand giebt als zwischen den Begriffen normativ und allgemeingültig. Die Moral auf „erste und letzte Erkenntnisse“ zu gründen, das musste desshalb misslingen, weil die Moral selbst nichts Erstes und Letztes, sondern etwas sehr Mittleres, Bedingtes, Relatives ist, nicht Ein kategorischer Imperativ, sondern eine Vielheit hypothetischer Imperative, ein System menschlicher Mittel zu menschlichen (keineswegs „absolut nothwendigen“) Zwecken: sei dieser Zweck nun die Erhaltung der ganzen menschlichen Gattung oder die Züchtung eines bestimmten Typus, eine chronische Vorkehrung zu Gunsten der Rasse oder eine Abwehr acuter Schädlichkeiten. Hier handelt es sich bestenfalls um richtige Deutung gesammelter Erfahrungen und angenäherte Vorausberechnung künftiger, aber nie und nimmer um etwas Endgültiges und für den „Menschen an sich“ oder gar für jedes vernünftige Wesen Verbindliches. Umgekehrt ist jede allgemeine Erkenntniss für alle besonderen Werthbestimmungsfragen, für das einfache Abschätzen und Gutheissen und Verbieten menschlicher Handlungen ohne Belang — es sei denn, man habe sie sich zuvor in der wünschenswerthen Weise verengert oder zurechtgemacht: mit bestochenen Zeugen lässt sich natürlich alles Mögliche beweisen. Eine theoretische Moralbegründung fristet ihr oft bändelanges Dasein damit, aus einer missdeuteten und präparirten Theorie eine bereits anderweitig anerkannte Moral abzuleiten; sie argumentirt\Foreign{ ex concessis,} und überdies falsch!
\par

